\chapter{Firmware ESP32}
\label{ch:04-firmware}

\section{Struktur Program}

Firmware adalah satu sketch Arduino dengan \code{setup()} dan \code{loop()}, tanpa
task FreeRTOS buatan sendiri. Semua pekerjaan berlangsung di dalam \code{loop()}
sebagai sederetan pemanggilan berurutan: satu langkah selesai lebih dahulu, baru
langkah berikutnya dijalankan. Setiap tick bersifat kooperatif --- ia memeriksa
\code{millis()} dan segera kembali bila periodenya belum tiba --- sehingga tidak ada
\code{delay()} di jalur utama dan tidak ada pekerjaan yang tertinggal berjalan di
latar belakang. Tumpukan jaringan ESP-IDF (WiFi, lwIP, TCP) memang berjalan pada task
di bawah sketch, tetapi sketch sendiri tidak membuat task maupun thread: dari sudut
pandang program hanya ada satu alur eksekusi.

\code{setup()} berjalan berurutan dan mencetak hasilnya ke serial pada
\SerialBaud{} baud: serial dibuka lalu diberi jeda singkat; perangkat keras
diinisialisasi (pindai bus I\textsuperscript{2}C, siapkan \code{TRIG} dan
\code{ECHO}, pasang servo pada posisi tertutup, deteksi DS3231, deteksi LCD pada
\AddrLcd{} lalu \AddrLcdAlt); konfigurasi dimuat dari NVS; jaringan disiapkan
dengan mode station, identitas dari MAC, nama host, pemindaian jaringan, alamat
broker, buffer \MqttBuffer{} byte, keepalive \KeepAliveS{} detik, socket timeout
sepuluh detik, dan callback pesan; splash ditampilkan pada LCD; lalu jarak diukur
sekali agar telemetri pertama tidak kosong.

\figref{fig:fw-flow-overview} memperlihatkan alur lengkap program: panel kiri adalah
\code{setup()} yang berjalan sekali saat boot, panel kanan adalah \code{loop()}. Kotak
bernomor pada panel kanan adalah urutan pemanggilan pada satu iterasi: panah selalu
bergerak dari langkah yang lebih kecil ke langkah yang lebih besar, lalu kembali ke
langkah 1, sehingga tidak ada dua langkah yang berjalan bersamaan. Dua deteksi tepi
naik ditangani di dalam langkah 1: saat koneksi MQTT baru terbentuk, telemetri langsung
dikirim sebagai penanda perangkat hidup; saat WiFi baru tersambung, permintaan
sinkronisasi waktu dimulai.

\figuretikz[\linewidth]{fw-flow-overview}{Alur lengkap firmware: \code{setup()} sekali
saat boot, lalu \code{loop()} satu thread dengan delapan pemanggilan berurutan pada
setiap iterasi. Kotak bernomor dapat memuat lebih dari satu pemanggilan berurutan;
rincian tiap fungsi ada di \tabref{tab:04-tick}.}{fig:fw-flow-overview}

\tabref{tab:04-tick} merinci setiap tick dan periodenya. Periode ``setiap iterasi''
berarti fungsi dipanggil pada setiap putaran \code{loop()} dan keluar sendiri bila
belum waktunya.

\begin{table}[H]
\centering
\caption{Tujuh tick kooperatif di dalam \code{loop()} beserta fungsinya.}
\label{tab:04-tick}
\small
\begin{tabular}{@{}p{3.1cm}p{2.4cm}p{9.1cm}@{}}
\toprule
\textbf{Fungsi} & \textbf{Periode} & \textbf{Tugas} \\
\midrule
\code{net::tick} & setiap iterasi & Menjaga koneksi: bila WiFi putus hanya
\code{client.loop()} yang dipanggil; bila MQTT putus, connect ulang paling cepat
\RetryMs{} ms. \\
\code{commandTick} & setiap iterasi & Mengambil satu perintah dari antrean
berukuran \CmdQueue{} lalu menjalankannya. \\
\code{ntpTick} & setiap iterasi & Menunggu hasil SNTP, menggeser UTC ke waktu
lokal, menulis hasilnya ke RTC, menampilkannya di LCD. \\
\code{feedTick} & setiap iterasi & Menjalankan state machine gate; hanya bekerja
saat fase tidak idle. \\
\code{sensorTick} & \SensorMs{} ms & Mengukur jarak dan menghitung persentase sisa
pakan. \\
\code{scheduleTick} & satu detik & Memeriksa jadwal lokal terhadap jam RTC dan
mengeksekusi paling banyak satu jadwal. \\
\code{telemetryTick} & \TelemetryMs{} ms & Mempublikasikan telemetri berkala. \\
\code{displayTick} & \DisplayMs{} ms & Menyegarkan LCD; halaman berganti tiap
\PageMs{} ms saat ada kondisi error. \\
\bottomrule
\end{tabular}
\end{table}

\src{firmware/esp32_feeder/esp32_feeder.ino}

\section{Mesin Keadaan Gate Pakan}

Gerakan gate dijalankan sebagai mesin keadaan berfase empat: idle, membuka,
menutup, dan jeda antar porsi. Nama fase pada kode mengikuti urutan itu, mulai
dari \code{FEED_IDLE} sampai \code{FEED_GAP}. Setiap fase hanya mencatat waktu
masuknya, dan \code{feedTick} memeriksa selisih \code{millis()} terhadap waktu
itu. Dengan pola ini gate dapat membuka dan menutup berkali-kali
tanpa menghentikan penerimaan perintah atau pembaruan LCD.
\figref{fig:fw-state-feed} menggambarkan transisinya dan
\tabref{tab:04-state} memerinci kondisinya.

\figureasset[\linewidth]{fw-state-feed}{Mesin keadaan gate pakan: fase idle, membuka, menutup,
dan jeda antar porsi.}{fig:fw-state-feed}

\begin{table}[H]
\centering
\caption{Transisi mesin keadaan gate pakan.}
\label{tab:04-state}
\small
\begin{tabular}{@{}p{2.4cm}p{2.4cm}p{3.9cm}p{5.5cm}@{}}
\toprule
\textbf{Dari} & \textbf{Ke} & \textbf{Kondisi} & \textbf{Aksi} \\
\midrule
\code{FEED_IDLE} & \code{FEED_OPEN} & \code{feedStart()} dipanggil, fase masih
idle, servo siap & Buka gate, catat waktu fase. \\
\code{FEED_OPEN} & \code{FEED_CLOSING} & Lewat \PortionOpenMs{} ms sejak gate
dibuka & Tutup gate, catat waktu fase. \\
\code{FEED_CLOSING} & \code{FEED_GAP} & Lewat \SettleMs{} ms dan porsi masih
tersisa & Catat waktu fase. \\
\code{FEED_CLOSING} & \code{FEED_IDLE} & Lewat \SettleMs{} ms dan porsi habis &
\code{finishFeed()}: tutup gate, ukur jarak, publikasikan hasil. \\
\code{FEED_GAP} & \code{FEED_OPEN} & Lewat \GapMs{} ms sejak gate ditutup & Buka
gate untuk porsi berikutnya. \\
\bottomrule
\end{tabular}
\end{table}

\figref{fig:fw-feed-timing} menempatkan transisi itu pada garis waktu satu perintah
dua porsi. Satu siklus porsi memerlukan \PortionCycleMs{} ms dan jeda \GapMs{} ms
disisipkan di antara dua porsi. Urutan publikasi pada akhir rangkaian sengaja
dibedakan menurut asal perintah: untuk pakan terjadwal, event dikirim lebih dahulu
baru telemetri, karena platform memakai baris telemetri terakhir sebagai level
sebelum; untuk pakan manual urutannya dibalik, agar saat ack tiba platform sudah
memiliki level sesudah.

\figuretikz[\linewidth]{fw-feed-timing}{Garis waktu satu perintah dua porsi: buka, tutup,
settle, jeda, lalu porsi kedua, diakhiri publikasi hasil dan ack.}{fig:fw-feed-timing}

\section{Konfigurasi Tersimpan}

Konfigurasi yang harus bertahan setelah daya putus disimpan di NVS pada namespace
\NvsNamespace, seperti pada \tabref{tab:04-nvs}. Nilai bawaan dipakai bila kunci
belum pernah ditulis, sehingga papan baru tetap berfungsi tanpa konfigurasi awal.

\begin{table}[H]
\centering
\caption{Kunci NVS pada namespace \code{feeder}.}
\label{tab:04-nvs}
\small
\begin{tabular}{@{}p{2.2cm}p{1.5cm}p{5.4cm}p{5.5cm}@{}}
\toprule
\textbf{Kunci} & \textbf{Tipe} & \textbf{Isi} & \textbf{Catatan} \\
\midrule
\code{threshold} & Int & Ambang persen pakan menipis & Bawaan \DefaultThreshold;
ditulis lewat \code{set_threshold}. \\
\code{fireDate} & UInt & Tanggal terakhir jadwal berjalan, dikodekan
\code{y*10000+mo*100+d} & Bersama bitmask mencegah jadwal berjalan dua kali sehari. \\
\code{fireMask} & UInt & Bitmask jadwal yang sudah dijalankan pada \code{fireDate}
& Satu bit untuk satu entri, maksimal \MaxSchedules{} entri. \\
\code{schedules} & String & Array JSON jadwal aktif & Ditimpa seluruhnya setiap
\code{set_schedule} diterima. \\
\bottomrule
\end{tabular}
\end{table}

\src{firmware/esp32_feeder/protocol.cpp}

Firmware juga memiliki anggaran memori untuk payload keluar. Buffer MQTT disetel
\MqttBuffer{} byte, sedangkan dokumen JSON telemetri, event, dan ack dibatasi
berturut-turut \ReserveTelemetry{}, \ReserveState{}, dan \ReserveAck{} byte.
Payload yang melebihi anggaran dibuang dan dicatat di log, bukan dipotong, karena
payload terpotong akan gagal di-parse platform. Anggaran ack paling besar karena
payload \code{get_state} memuat snapshot keadaan lengkap di dalam field
\code{message}.

\section{Penjadwalan Lokal}

Jadwal harian dijalankan perangkat, bukan platform. Platform hanya menyimpan jadwal
dan mengirimkannya lewat \code{set_schedule}; setelah itu perangkat memegang kendali
waktu. Pilihan ini membuat jadwal tetap berjalan saat jaringan atau host mati, dan
menghilangkan kebutuhan platform terjaga sepanjang hari.

Setiap detik \code{scheduleTick} membaca jam RTC; bila jam belum valid, fungsi
langsung kembali sehingga tidak ada jadwal dieksekusi dengan waktu salah. Untuk
setiap entri, jam dan menit harus sama dengan jam sekarang dan detik harus masih di
dalam jendela \WindowS{} detik pertama menit itu, sehingga jadwal 07:00 dapat
berjalan antara 07:00:00 dan 07:00:59. Jumlah entri dibatasi \MaxSchedules{}, dan
setiap entri menempati satu bit pada bitmask harian sehingga tidak dapat berjalan
dua kali pada hari yang sama meskipun perangkat di-reboot.

Empat keadaan membuat jadwal tidak dieksekusi. Bila gate sibuk, eksekusi ditunda ke
detik berikutnya selama masih di dalam jendela. Bila servo tidak terpasang atau
hopper terbaca kosong, jadwal ditandai sudah dijalankan hari itu dan dilewati,
karena mencoba lagi hanya mengulang kegagalan yang sama. Bila jam belum valid, tidak
ada jadwal yang dicoba. Pembedaan antara ``ditunda'' dan ``ditandai dijalankan''
disengaja: penundaan untuk keadaan yang dapat berubah dalam hitungan detik,
penandaan untuk keadaan yang tidak berubah tanpa tindakan manusia.
\figref{fig:fw-seq-scheduled-feed} memperlihatkan akibatnya di sisi platform.

\figureasset[\linewidth]{fw-seq-scheduled-feed}{Urutan eksekusi jadwal lokal: perangkat
menjalankan gate atas dasar jam RTC, lalu melaporkan event ke topik state.}{fig:fw-seq-scheduled-feed}

\section{Siklus Hidup Koneksi MQTT}

\figref{fig:fw-seq-mqtt-lifecycle} merangkum siklus hidup koneksi. Tiga hal
ditetapkan saat connect: LWT berupa string \code{offline} pada topik
\topic{status} dengan tanda retained; publikasi birth \code{online} pada topik yang
sama, juga retained; dan langganan pada topik \topic{command}. Flag retained pada
keduanya membuat platform selalu memiliki gambaran terakhir, karena subscriber baru
langsung menerima \code{online} atau \code{offline}. Percobaan connect diulang
paling cepat setiap \RetryMs{} ms, dan keepalive \KeepAliveS{} detik membuat broker
menganggap sesi mati bila tidak ada paket sama sekali pada rentang itu.

\figureasset[\linewidth]{fw-seq-mqtt-lifecycle}{Siklus hidup koneksi MQTT perangkat: connect
dengan LWT terdaftar, publikasi birth retained, langganan perintah, lalu reconnect
berkala setelah koneksi putus.}{fig:fw-seq-mqtt-lifecycle}

Satu keputusan perlu dijelaskan: tidak ada pemutusan koneksi rapi ketika WiFi
hilang. Firmware sengaja tidak memanggil \code{disconnect()} pada keadaan itu,
karena paket \code{DISCONNECT} membatalkan wasiat: perangkat akan terlihat
memutuskan diri dengan sengaja, LWT tidak diterbitkan, dan pesan retained
\code{online} tertinggal sehingga dashboard tidak pernah melihat perangkat jatuh.
Yang dilakukan firmware adalah membiarkan sesi TCP mati sendiri, membiarkan broker
menerbitkan \code{offline}, dan tetap memanggil \code{client.loop()} agar pustaka
MQTT menyadari socket sudah mati lalu menyambung ulang.
(\file{firmware/esp32_feeder/net.cpp})

Callback pesan menyaring sebelum apa pun terjadi: pesan yang topiknya bukan
\code{command} langsung dibuang. Penyaringan itu menangkal perilaku broker
tertanam yang dapat mengirim pesan retained dari topik lain kepada subscriber mana
pun. Setelah lolos saringan, payload masuk ke antrean ring berukuran \CmdQueue{}
dan dijalankan satu per satu oleh \code{commandTick}; bila antrean penuh, perintah
baru dibuang dan dicatat di log.

\section{Mode Degradasi}

Setiap periferal bersifat opsional. Dengan hanya papan pengendali terpasang,
firmware tetap menyambung, tetap menjawab setiap perintah, dan tidak pernah
reboot; yang hilang hanya kemampuan terkait (\tabref{tab:04-degradasi}).

\begin{table}[H]
\centering
\caption{Perilaku firmware ketika sebuah periferal tidak terdeteksi.}
\label{tab:04-degradasi}
\small
\begin{tabular}{@{}p{2.4cm}p{12.6cm}@{}}
\toprule
\textbf{Periferal hilang} & \textbf{Akibat} \\
\midrule
Sensor ultrasonik & Telemetri tetap dikirim tanpa \code{distance_cm} dan berstatus
\code{error}; level pakan tidak diketahui; pakan manual tetap diizinkan karena tidak
ada pembacaan level yang dapat dipakai sebagai penjaga. \\
Motor servo & Pakan manual dijawab \code{error_servo}; jadwal dilewati dan ditandai
sudah dijalankan untuk hari itu. \\
Modul RTC & \code{set_rtc} tetap berhasil dengan jam internal di RAM, tetapi jam
hilang saat reboot; jadwal baru berjalan setelah \code{set_rtc} atau sinkronisasi
NTP. \\
LCD karakter & Tidak ada tampilan lokal; seluruh fungsi lain berjalan biasa. \\
\bottomrule
\end{tabular}
\end{table}

\section{Pustaka dan Versi}

Firmware dibangun dengan inti Arduino esp32 \CoreVersion{} dan opsi papan
\code{\Fqbn}. \tabref{tab:04-pustaka} memuat pustaka pihak ketiga yang dipakai.

\begin{table}[H]
\centering
\caption{Pustaka firmware dan versinya.}
\label{tab:04-pustaka}
\small
\begin{tabular}{@{}p{4.4cm}p{4.6cm}p{1.6cm}p{3.7cm}@{}}
\toprule
\textbf{Pustaka} & \textbf{Penulis} & \textbf{Versi} & \textbf{Kegunaan} \\
\midrule
PubSubClient & Nick O'Leary & 2.8+ & Klien MQTT sinkron, LWT, buffer. \\
ArduinoJson & Benoit Blanchon & 7.x & Parse perintah dan susun payload. \\
RTClib & Adafruit & 2.1+ & Akses DS3231. \\
LiquidCrystal I2C & Frank de Brabander & 1.1+ & LCD1602 melalui PCF8574. \\
ESP32Servo & K. Harrington, J. K. Bennett & 0.13+ & PWM servo berbasis LEDC. \\
\bottomrule
\end{tabular}
\end{table}

\src{firmware/esp32_feeder/README.md}
